\documentclass[11pt,a4paper]{article}

\usepackage[utf8]{inputenc}
\usepackage[T1]{fontenc}
\usepackage[spanish,es-tabla,es-nodecimaldot]{babel}
\usepackage{lmodern}
\usepackage[margin=2.5cm]{geometry}
\usepackage{amsmath,amssymb}
\usepackage{booktabs}
\usepackage{siunitx}
\usepackage[hidelinks]{hyperref}

\sisetup{output-decimal-marker={,}, group-separator={\,}, group-minimum-digits=4, group-digits=integer}

\newcommand{\WO}{W_0}

\title{Dimensionamiento preliminar de aeronaves de transporte\\[4pt]
       \large Metodología implementada en el código \texttt{sizing}}
\author{Diseño, Cálculo y Certificación de Aeronaves y Vehículos Espaciales}
\date{}

\begin{document}
\maketitle

\begin{abstract}
Este documento resume las ecuaciones que implementa el código de
dimensionamiento. El método es el de Raymer~\cite{raymer6,raymer7} en sus dos
niveles: el \emph{first-order sizing} (capítulo~3), que obtiene el peso máximo
al despegue $\WO$ a partir de la misión y de correlaciones estadísticas, y el
\emph{initial sizing} (capítulos~5 y~6), que fija la carga alar y la relación
empuje-peso, refina las fracciones de peso y dimensiona ala, fuselaje y
empenaje. Donde las ediciones~6 y~7 del libro difieren se dan ambos juegos de
coeficientes; el caso de estudio de referencia es el del guion de la
asignatura~\cite{guion}.
\end{abstract}

\section{Planteamiento}

El peso al despegue se descompone en tripulación, carga de pago, combustible y
peso en vacío. Expresando los dos últimos como fracciones de $\WO$
\cite[ec.~3.4]{raymer6}:
\begin{equation}
  \WO = \frac{W_\text{crew} + W_\text{payload}}
             {1 - \left(W_f/\WO\right) - \left(W_e/\WO\right)} .
  \label{eq:w0}
\end{equation}
Como $W_e/\WO$ depende de $\WO$, la ecuación se resuelve por punto fijo: se
supone $\WO$, se evalúa el segundo miembro y se repite hasta que la variación
es menor que la tolerancia (\SI{0.01}{kg} en el código). Los dos niveles del
método solo difieren en cómo se estiman $W_f/\WO$ y $W_e/\WO$
(tabla~\ref{tab:niveles}).

\begin{table}[h]
\centering
\caption{Diferencias entre los dos niveles de cálculo.}
\label{tab:niveles}
\begin{tabular}{lll}
\toprule
Término & First-order sizing & Initial sizing \\
\midrule
Ascenso            & \num{0.985} (tabla 3.2)            & $1{,}0065-0{,}0325M$ (ec.~6.9) \\
$(L/D)$ de crucero & $0{,}866\,(L/D)_\text{máx}$        & Polar parabólica (ec.~6.13) \\
$W_e/\WO$          & Tabla 3.1 (solo $\WO$)             & Tabla 6.1 ($\WO$, $A$, $T/W$, $W/S$, $M_\text{máx}$) \\
\bottomrule
\end{tabular}
\end{table}

Todas las masas se expresan en \si{kg} y la carga alar en \si{kg/m^2}; se
multiplica por $g=\SI{9.81}{m/s^2}$ allí donde la ecuación exige fuerza por
unidad de superficie. Las propiedades del aire (densidad y velocidad del
sonido a la altitud de crucero) se toman de la atmósfera estándar
internacional~\cite{isa}.

%----------------------------------------------------------------------------
\section{First-order sizing}

\subsection{Tripulación y carga de pago}
\begin{equation}
  W_\text{payload} = n_\text{pax}\, m_\text{pax}, \qquad
  W_\text{crew} = n_\text{crew}\, m_\text{crew}.
\end{equation}
El guion~\cite{guion} adopta \SI{100}{kg} por pasajero con su equipaje y
\SI{95}{kg} por tripulante; el número mínimo de tripulantes de cabina es uno
por cada 50 asientos instalados~\cite{airops}.

\subsection{Fracciones de peso de la misión}
Cada segmento $i$ se representa por la fracción $W_i/W_{i-1}$; su producto da
la fracción al final de la misión, $W_x/\WO$. Los segmentos cortos toman los
valores históricos de la tabla~3.2 de Raymer~\cite{raymer6}: arranque y
despegue \num{0.970}, ascenso \num{0.985} y aterrizaje \num{0.995}. Raymer
incluye el descenso en el crucero; el código permite además un segmento de
descenso explícito (\num{0.990}), desactivable por caso.

\paragraph{Crucero.} Ecuación de Breguet del alcance para un reactor
\cite[ec.~3.6]{raymer6}:
\begin{equation}
  \frac{W_\text{cr}}{W_\text{cr-1}} = \exp\!\left[\frac{-R\,C}{V\,(L/D)}\right],
  \qquad V = M\,a(h),
\end{equation}
con $R$ el alcance, $C$ el consumo específico de empuje (tabla~3.3,
\SI{0.5}{h^{-1}} en crucero para un turbofán de alto índice de derivación) y
$a(h)$ la velocidad del sonido ISA a la altitud de crucero.

\paragraph{Espera.} Ecuación de la autonomía \cite[ec.~3.8]{raymer6}, con
$E$ el tiempo de espera y $C=\SI{0.4}{h^{-1}}$ (tabla~3.3):
\begin{equation}
  \frac{W_\text{lt}}{W_\text{lt-1}} = \exp\!\left[\frac{-E\,C}{(L/D)}\right].
\end{equation}

\subsection{Eficiencia aerodinámica}
La $(L/D)$ máxima se estima con el alargamiento mojado
\cite[ec.~3.12]{raymer6}:
\begin{equation}
  (L/D)_\text{máx} = K_{LD}\sqrt{\frac{A}{S_\text{wet}/S_\text{ref}}},
  \qquad K_{LD} = 15{,}5 \text{ (reactor civil)},
\end{equation}
con $S_\text{wet}/S_\text{ref}$ de la figura~3.6. Para un reactor el crucero
de máximo alcance se vuela a $0{,}866\,(L/D)_\text{máx}$ y la espera a
$(L/D)_\text{máx}$ \cite[cap.~3]{raymer6}.

\subsection{Fracción de combustible}
Con un \SI{6}{\percent} para reservas y combustible atrapado
\cite[ec.~3.13]{raymer6}:
\begin{equation}
  \frac{W_f}{\WO} = 1{,}06\left(1 - \frac{W_x}{\WO}\right),
  \qquad \frac{W_x}{\WO} = \prod_i \frac{W_i}{W_{i-1}} .
\end{equation}

\subsection{Fracción de peso en vacío}
Correlación estadística de la tabla~3.1 para transporte a reacción, en
unidades métricas:
\begin{equation}
  \frac{W_e}{\WO} = a\,\WO^{\,C}\,K_{vs}\,K_c,
\end{equation}
con $a=\num{0.97}$, $C=\num{-0.06}$ en la 6.ª edición~\cite{raymer6} y
$a=\num{1.202}$, $C=\num{-0.072}$ en la 7.ª~\cite{raymer7}. $K_{vs}=1{,}04$
para ala de flecha variable y 1 para ala fija. $K_c=0{,}95$ aproxima una
estructura de materiales compuestos y $K_c=1$ una metálica
\cite[cap.~3]{raymer6}; el mismo factor se aplica a la fracción refinada de la
tabla~6.1. Con ambas fracciones se itera la ecuación~\eqref{eq:w0}.

%----------------------------------------------------------------------------
\section{Initial sizing}

\subsection{Carga alar}

\paragraph{Aterrizaje.} La velocidad de aproximación limita la carga alar. De
$L=W$ en pérdida con $C_{L,\text{máx}}$ de aterrizaje \cite[ec.~5.6]{raymer6}:
\begin{equation}
  \left(\frac{W}{S}\right)_\text{L} = \frac{1}{g}\,\frac{1}{2}\rho_0\,
  V_\text{stall}^2\,C_{L,\text{máx}},
  \qquad V_\text{stall} = \frac{V_\text{app}}{1{,}23},
\end{equation}
con $\rho_0$ al nivel del mar y el factor $1{,}23$ de CS-25~\cite{cs25}. El
valor se refiere al despegue dividiendo por $W_\text{MLW}/\WO$.

\paragraph{Crucero.} La resistencia se modela con la polar parabólica
$C_D = C_{D0} + K C_L^2$, donde
\begin{equation}
  C_{D0} = C_{fe}\,\frac{S_\text{wet}}{S_\text{ref}}, \qquad
  K = \frac{1}{\pi A e},
\end{equation}
con $C_{fe}=\num{0.0026}$ para transporte civil \cite[ec.~12.23,
tabla~12.3]{raymer6}. El factor de Oswald sigue las correlaciones de Raymer
\cite[ecs.~12.48--12.49]{raymer6}:
\begin{equation}
  e =
  \begin{cases}
    1{,}78\left(1-0{,}045A^{0{,}68}\right) - 0{,}64,
      & \Lambda_\text{LE} \le 30^\circ, \\[4pt]
    4{,}61\left(1-0{,}045A^{0{,}68}\right)\left(\cos\Lambda_\text{LE}\right)^{0{,}15} - 3{,}1,
      & \Lambda_\text{LE} > 30^\circ,
  \end{cases}
\end{equation}
donde la flecha del borde de ataque se obtiene de la de la línea del
\SI{25}{\percent} para un ala trapezoidal:
\begin{equation}
  \tan\Lambda_\text{LE} = \tan\Lambda_{c/4} + \frac{1-\lambda}{A\,(1+\lambda)} .
\end{equation}
Raymer indica que $e$ suele estar entre \num{0.70} y \num{0.85}
\cite[sec.~12.6.1]{raymer6}; el código avisa cuando la correlación da un valor
fuera de ese rango. Entre \SIlist{0;30}{\degree} de flecha el libro propone
interpolar entre las dos ecuaciones; no se hace porque en alas como la del guion
($\Lambda_\text{LE}\approx\SI{28}{\degree}$) daría $e\approx0{,}49$, fuera de ese
mismo rango, y el guion usa la de ala recta~\cite{guion}.
Cuando la correlación cae fuera del rango, el caso puede fijar $e$ a mano
(\texttt{Aerodynamics.oswald}). Es lo que hace el ejemplo del A380: con
$\Lambda_\text{LE} > \SI{30}{\degree}$ y $A = 7{,}53$ la ecuación~12.49 da
$e \approx 0{,}57$, y el caso toma $e = 0{,}80$, el centro de la banda típica.
La carga alar de máximo alcance para un reactor corresponde a
$C_{L,\text{opt}} = \sqrt{C_{D0}/(3K)}$ \cite[ec.~5.14]{raymer6}:
\begin{equation}
  \left(\frac{W}{S}\right)_\text{cr} = \frac{q}{g}\,\sqrt{\frac{\pi A e\,C_{D0}}{3}},
  \qquad q = \tfrac{1}{2}\rho V^2,
\end{equation}
que se refiere al despegue dividiendo por el producto de las fracciones de
despegue y ascenso.

\paragraph{Selección.} La carga alar de diseño se elige por debajo del límite
de aterrizaje. El crucero pediría un ala mayor; se acepta volar a un $C_L$
algo superior al óptimo a cambio de un ala más ligera y de menor
resistencia~\cite{guion}.

\subsection{Relación empuje-peso}
Correlación estadística frente al Mach máximo (tabla~5.3):
\begin{equation}
  \frac{T}{\WO} = a\,M_\text{máx}^{\,C},
\end{equation}
con $a=\num{0.267}$, $C=\num{0.363}$ en la 6.ª edición y $a=\num{0.297}$,
$C=\num{0.350}$ en la 7.ª. En un birreactor la subida con un motor inoperativo
exigida por CS-25~\cite{cs25} eleva el valor adoptado por encima del
estadístico; el empuje total es $T = (T/W)\,\WO\,g$.

\subsection{Geometría del ala}
Para un ala trapezoidal de alargamiento $A$ y estrechamiento $\lambda$:
\begin{equation}
  S = \frac{\WO}{W/S}, \quad
  b = \sqrt{A\,S}, \quad
  c_r = \frac{2S}{b\,(1+\lambda)}, \quad
  c_t = \lambda\,c_r, \quad
  \bar{c} = \frac{2}{3}\,c_r\,\frac{1+\lambda+\lambda^2}{1+\lambda}.
\end{equation}

\subsection{Fracción de peso en vacío refinada}
La tabla~6.1 añade la influencia del alargamiento, la relación empuje-peso,
la carga alar y el Mach máximo. En la 6.ª edición~\cite{raymer6}:
\begin{equation}
  \frac{W_e}{\WO} = \left[a + b\,\WO^{\,C_1} A^{C_2}
    \left(\frac{T}{\WO}\right)^{C_3}\left(\frac{\WO}{S}\right)^{C_4}
    M_\text{máx}^{\,C_5}\right] K_{vs}\,K_c,
\end{equation}
con $a=\num{0.32}$, $b=\num{0.66}$, $C_1=\num{-0.13}$, $C_2=\num{0.30}$,
$C_3=\num{0.06}$, $C_4=\num{-0.05}$, $C_5=\num{0.05}$. La tabla está en
unidades imperiales ($\WO$ en \si{lb}, $\WO/S$ en \si{lb/ft^2}); con la
conversión de más abajo, en \si{kg} y \si{kg/m^2} queda $b_\text{SI} =
\num{0.645}$. La 7.ª edición~\cite{raymer7} elimina el término aditivo:
\begin{equation}
  \frac{W_e}{\WO} = a\,\WO^{\,C_1} A^{C_2}
    \left(\frac{T}{\WO}\right)^{C_3}\left(\frac{\WO}{S}\right)^{C_4}
    M_\text{máx}^{\,C_5}\,K_{vs}\,K_c,
\end{equation}
con $a=\num{0.869}$, $C_1=\num{-0.037}$, $C_2=\num{0.398}$,
$C_3=\num{0.100}$, $C_4=\num{-0.161}$, $C_5=\num{0.050}$. Estos coeficientes
solo se dan en unidades imperiales ($\WO$ en \si{lb}, $\WO/S$ en
\si{lb/ft^2}); en las dos ediciones, para usar \si{kg} y \si{kg/m^2} se
convierte la constante que multiplica a $\WO^{\,C_1}(\WO/S)^{C_4}$ ($b$ en la
6.ª, $a$ en la 7.ª):
\begin{equation}
  k_\text{SI} = k_\text{fps}\,k_m^{\,C_1}\left(\frac{k_m}{k_S}\right)^{C_4},
  \qquad k_m = \SI{2.20462}{lb/kg}, \quad k_S = \SI{10.7639}{ft^2/m^2},
\end{equation}
lo que da $a_\text{SI} = \num{1.089}$ en la 7.ª edición.

\subsection{Fracciones de misión refinadas}

\paragraph{Despegue.} Se mantiene \num{0.970}, dentro del rango
$0{,}97$--$0{,}99$ de la ecuación~6.8.

\paragraph{Ascenso y aceleración.} Hasta el Mach de crucero, partiendo de
$M=0{,}1$ \cite[ec.~6.9]{raymer6}:
\begin{equation}
  \frac{W_\text{cl}}{W_\text{cl-1}} = 1{,}0065 - 0{,}0325\,M .
\end{equation}

\paragraph{Crucero.} Misma ecuación de Breguet \cite[ec.~6.11]{raymer6},
pero con la $L/D$ evaluada con la polar en la condición de vuelo, usando
$L=W$ \cite[ec.~6.13]{raymer6}:
\begin{equation}
  \frac{L}{D} = \frac{1}{\dfrac{q\,C_{D0}}{W/S} + \dfrac{W}{S}\,\dfrac{1}{q\,\pi A e}} .
\end{equation}
La carga alar de esta expresión es la real en la condición evaluada, no la del
despegue \cite[cap.~6]{raymer6}: la de despegue multiplicada por las fracciones
de los segmentos anteriores \cite[cap.~5]{raymer6}. Se evalúa en un único punto
a mitad de crucero, como admite Raymer para el análisis inicial
\cite[cap.~12]{raymer6}:
\begin{equation}
  \left(\frac{W}{S}\right)_\text{cr} = \frac{\WO}{S}\,
  \frac{W_\text{to}}{\WO}\,\frac{W_\text{cl}}{W_\text{cl-1}}\,
  \frac{1 + W_\text{cr}/W_\text{cr-1}}{2},
\end{equation}
que se itera porque la fracción de crucero depende de la $L/D$. El guion fija
en su lugar \SI{550}{kg/m^2}~\cite{guion}; el código permite forzar ese valor.

\paragraph{Espera.} Misma ecuación de la autonomía \cite[ec.~6.14]{raymer6},
volando a la $(L/D)$ máxima \cite[cap.~5]{raymer6}, que ahora se toma de la
polar para ser coherente con el crucero:
\begin{equation}
  \left(\frac{L}{D}\right)_\text{máx} = \frac{1}{2\sqrt{C_{D0}\,K}} .
\end{equation}

\paragraph{Aterrizaje y reservas.} Sin cambios respecto al first-order sizing. Con las nuevas fracciones se recalcula $W_f/\WO$ y se
vuelve a iterar la ecuación~\eqref{eq:w0} con la fracción en vacío de la
tabla~6.1; con $\WO$ refinado se recalculan la superficie alar, la geometría
del ala y el empuje.

\subsection{Fuselaje}
En un avión de pasajeros el fuselaje lo define la cabina; la estadística solo
sirve de comprobación~\cite{guion}.

\paragraph{Sección.} El ancho interior de cabina es la suma de asientos,
pasillos y holguras; el exterior añade estructura y aislamiento. La altura
apila quilla, bodega (contenedor LD3), suelo y cabina de cada cubierta y
techo. Si la sección no es circular se usa un diámetro equivalente obtenido
del área de la sección \cite[cap.~6]{raymer6}:
\begin{equation}
  d_\text{eq} = \sqrt{w\,h} .
\end{equation}

\paragraph{Longitud.} La longitud de cabina suma las filas de cada clase
($n_\text{filas}\times$ paso), cocinas, aseos, pares de salidas tipo~A
(uno por cada 110 asientos, CS~25.807~\cite{cs25}) y escaleras. La longitud
total añade morro y cono de cola:
\begin{equation}
  L_f = L_\text{cabina} + L_\text{morro} + L_\text{cono}.
\end{equation}
El morro y el cono se obtienen del criterio de contorno de Raymer
\cite[fig.~8.2]{raymer6} (ángulos máximos de las superficies superior e inferior respecto a la corriente) o
se fijan como decisión de proyecto.

\paragraph{Comprobaciones.} Relación de finura $L_f/d$, cuyo óptimo
aerodinámico subsónico está entre 6 y 8 \cite[cap.~6]{raymer6}, y longitud
estadística de la tabla~6.3:
\begin{equation}
  L_f = a\,\WO^{\,C},
\end{equation}
con $a=\num{0.287}$, $C=\num{0.43}$ en la 6.ª edición y $a=\num{0.69}$,
$C=\num{0.36}$ en la 7.ª (\si{m}, \si{kg}).

\subsection{Empenaje}
Método de los coeficientes de volumen de cola
\cite[ecs.~6.28--6.29]{raymer6}:
\begin{equation}
  S_\text{HT} = \frac{c_\text{HT}\,\bar{c}\,S}{L_\text{HT}}, \qquad
  S_\text{VT} = \frac{c_\text{VT}\,b\,S}{L_\text{VT}},
\end{equation}
con $c_\text{HT}=1{,}00$ y $c_\text{VT}=0{,}09$ para transporte a reacción
(tabla~6.4). El brazo de cola se toma como fracción de la longitud del
fuselaje: \SIrange{50}{55}{\percent} con motores en el ala y
\SIrange{45}{50}{\percent} con motores traseros \cite[cap.~6]{raymer6}.
Por defecto la longitud es la obtenida por disposición de cabina; opcionalmente
puede tomarse la estadística de la tabla~6.3, $L_f = a\,W_0^{C_1}$, evaluada
con el $W_0$ de diseño.

Los coeficientes de la tabla~6.4 corresponden a una cola convencional. Raymer
\cite[sec.~6.4]{raymer6} permite reducirlos según la configuración: una cola en
T reduce un \SI{5}{\percent} ambos coeficientes (efecto placa del
estabilizador horizontal sobre el vertical y aire limpio en el horizontal),
una cola en H reduce un \SI{5}{\percent} el horizontal, un estabilizador
horizontal todo móvil lo reduce entre un 10 y un \SI{15}{\percent} (el código
toma \SI{12.5}{\percent}) y un sistema de mandos de vuelo activo permite
reducir ambas superficies en torno a un \SI{10}{\percent}. Las reducciones se
aplican multiplicando. Una cola en V se dimensiona como una convencional y se
sustituye por dos superficies con la misma área total, $S_\text{HT} +
S_\text{VT}$, y un diedro $\Gamma = \arctan\sqrt{S_\text{VT}/S_\text{HT}}$,
normalmente próximo a \SI{45}{\degree}.

Las superficies de mando mantienen una cuerda relativa constante
\cite[sec.~6.6]{raymer6}. Los timones ocupan toda la envergadura de la cola, así
que su fracción de área es la de cuerda de la tabla~6.5: \num{0.25} para el de
profundidad y \num{0.32} para el de dirección en transporte a reacción. Los
alerones van del \SIrange{50}{90}{\percent} de la semienvergadura, con
$c_a/c \approx 0{,}23$ (figura~6.3); para el ala trapezoidal,
\begin{equation}
  S_a = \frac{c_a}{c}\,S\,
  \frac{\left[\eta - (1-\lambda)\,\eta^2/2\right]_{0{,}5}^{0{,}9}}{1-(1-\lambda)/2} .
\end{equation}
Como alternativa, los alerones pueden fijarse directamente como una fracción
de $S$ (\texttt{aileron\_area\_fraction}); el guion V4 toma
$S_a = 0{,}05\,S$.

%----------------------------------------------------------------------------
\section{Procedimiento}
\begin{enumerate}
  \item Requisitos de misión: pasajeros, alcance, Mach, altitud, espera,
        velocidad de aproximación.
  \item First-order sizing: $W_\text{crew}+W_\text{payload}$, fracciones de
        misión, $W_f/\WO$, tabla~3.1 e iteración de~\eqref{eq:w0}.
  \item Initial sizing: límites de carga alar, elección de $W/S$ y $T/W$,
        tabla~6.1, fracciones refinadas e iteración de~\eqref{eq:w0}.
  \item Geometría con el $\WO$ refinado: ala, fuselaje y empenaje, en este
        orden, porque el empenaje necesita la cuerda media y la envergadura
        del ala y la longitud del fuselaje.
\end{enumerate}

%----------------------------------------------------------------------------
\begin{thebibliography}{9}

\bibitem{raymer6}
D.~P. Raymer.
\newblock \emph{Aircraft Design: A Conceptual Approach}, 6.ª ed.
\newblock AIAA Education Series. American Institute of Aeronautics and
  Astronautics, Reston, VA, 2018.

\bibitem{raymer7}
D.~P. Raymer.
\newblock \emph{Aircraft Design: A Conceptual Approach}, 7.ª ed.
\newblock AIAA Education Series. American Institute of Aeronautics and
  Astronautics, Reston, VA.

\bibitem{guion}
\emph{Caso completo de dimensionamiento de una aeronave de transporte de
  pasajeros} (versión~4).
\newblock Diseño, Cálculo y Certificación de Aeronaves y Vehículos
  Espaciales, Máster Universitario en Ingeniería Aeronáutica.

\bibitem{isa}
ISO 2533:1975.
\newblock \emph{Standard Atmosphere}.
\newblock International Organization for Standardization, 1975.

\bibitem{cs25}
EASA.
\newblock \emph{Certification Specifications and Acceptable Means of
  Compliance for Large Aeroplanes (CS-25)}.
\newblock European Union Aviation Safety Agency.

\bibitem{airops}
Reglamento (UE) n.º 965/2012 de la Comisión (Air OPS), Anexo~III (Part-ORO),
ORO.CC.100, y Anexo~IV (Part-CAT), CAT.POL.MAB.100.

\end{thebibliography}

\end{document}
